\begin{pastexam}[label={px:ncu-114-dsa-09}]{中央資工所 114 學年度 第 9 題}
（複選；每題 5 分，全部答對才給分，答錯倒扣 1 分）

（第 8--10 題共用下面的說明和程式。）

Consider a UAV (Unmanned Aerial Vehicle) with F unit of fuel travels from NCU to a destination
``target'' km away. There are n gas stations along the way. When the UAV refuels at a gas station, all
the fuel of the gas station is transferred into the UAV.

Suppose that the UAV consumes one unit of fuel for every kilometer it travels. Let the position and fuel
of a gas station indicate the distance between the gas station and NCU and its volume of fuel,
respectively. Return the minimum number of refueling stops the UAV must make in order to reach its
destination.
\begin{minted}[fontsize=\small]{c}
typedef struct {
    int position;
    int fuel;
} stop_info;

int Refuel(int target, int F, int n, stop_info s[]) {
    int N_R = 0;   // number of refuel
    int i;   //gas stations ID
    stop_info X;

    while (F < target) {
        for (i=0;  L1  ; ++i)
            Q.push(s[i].fuel);
        if (Q.empty()) return -1;
        X=Q.pop();
        F += X.fuel
        N_R++;
    }
    return N_R;
}
\end{minted}

Consider Refuel algorithm above. Which of the following statements are true?
\begin{enumerate}[label=(\alph*)]
\item Q could be a stack
\item Q could be a priority queue
\item Q could be a max heap
\item Q could be min heap
\end{enumerate}
\end{pastexam}

\begin{pastsoln}
答案：(b)(c)。

程式的讀法見歷屆試題~\ref{px:ncu-114-dsa-08} 的詳解：\texttt{L1} 填 \texttt{(i<n) \&\& (s[i].position <=F)}，
每次油不夠時，從已經經過的站裡挑\textbf{油最多}的那一站加油。 所以 \texttt{Q.pop()} 必須傳回最大值，
\texttt{Q} 要是最大優先佇列（資結書第~8.3~節）。

\begin{itemize}
\item (a) 堆疊彈出的是最後推入的，也就是最遠的那一站，不一定油最多。 反例：$\mathit{target} = 100$、
$F = 50$，兩站 $(\mathit{position}, \mathit{fuel}) = (10, 50)$、$(40, 10)$。 兩站都到得了，依序推入
$50$、$10$。 最佳解只在第一站加油，$F = 100$，$1$ 次。 用堆疊時先彈出 $10$，$F = 60$ 不夠，再彈出
$50$，共 $2$ 次。 錯。
\item (b) 優先佇列以油量為優先權、每次取最大的，正是需要的操作。 對。
\item (c) 最大堆積是最大優先佇列的標準實作，\texttt{pop} 取出最大值，$O(\log n)$。 對。
\item (d) 最小堆積每次取油最少的站，同一個反例會先取 $10$，共 $2$ 次。 錯。
\end{itemize}
(b) 的「priority queue」沒有說是最大還是最小；這裡讀成以油量大的優先，所以算對。
\end{pastsoln}
